% ============================================================
% SCHEDE PG AGGIUNTIVE — CASA SPINOLA
% Personaggi per tavolo alternativo / quinta sessione
% ============================================================

\chapter{Schede PG Aggiuntive --- Casa Spinola}
\index{schede PG!Spinola!aggiuntivi}

% ════════════════════════════════════════════════════════════
% PG 6 — RANIERI COSTA detto IL LUNGO
% @rpg-schema: <https://rpg-schema.org/instances/genova1507/pg-ranieri-costa>
%   a fitd:Character ;
%   rdfs:label "Ranieri Costa detto il Lungo" ;
%   fitd:crew <.../crew-spinola> ;
%   fitd:role "Il Bravo / Il Muscolo degli Spinola" ;
%   fitd:action [ fitd:actionName "Skirmish" ; fitd:dots 4 ] ;
%   fitd:action [ fitd:actionName "Prowl" ; fitd:dots 3 ] ;
%   fitd:action [ fitd:actionName "Hunt" ; fitd:dots 2 ] .
% ════════════════════════════════════════════════════════════

\pgheader[SpinolaRed]{Ranieri Costa detto il Lungo}{Il Bravo / Il Muscolo degli Spinola}{Spinola}
\pgquote{Raffaele Spinola risolve i problemi con il denaro. Io risolvo i problemi che il denaro non può risolvere.}

\section*{Background \& Motivazione}

\begin{rulebox}{}
  \textbf{Background:} Trentadue anni, ligure di Finale, figlio di un mercenario.
  È al servizio degli Spinola da otto anni — non come guardia del corpo ufficiale,
  ma come «consulente per questioni delicate». Questo significa: recupero crediti
  che nessun notaio può gestire, messaggi consegnati a persone che preferirebbero
  non essere trovate, e — raramente ma succede — la risoluzione permanente di
  problemi particolarmente ostinati. Non è un assassino di professione.
  È un professionista della violenza calibrata.

  \medskip\noindent\textcolor{SpinolaRed}{\textbf{Motivazione:}} La rivolta gli
  offre qualcosa che il lavoro ordinario non dà: la possibilità di usare
  le sue competenze per qualcosa di più grande di un recupero crediti.
  Non è idealismo — è noia. Otto anni di lavori ordinari lo hanno stancato.
\end{rulebox}

\section*{Attributi \& Azioni}


\begin{attributoblock}[SpinolaRed]{INSIGHT — Mente}
  \azione{Caccia}{Hunt}{2}{Trova persone, segue tracce, localizza in città}
  \azione{Studio}{Study}{1}{Sa leggere, conta i soldi, niente di più}
  \azione{Osserva}{Survey}{2}{Valuta minacce e uscite di emergenza istintivamente}
  \azione{Ingegno}{Tinker}{1}{Trappole semplici, meccanismi di serratura}
\end{attributoblock}

\columnbreak

\begin{attributoblock}[SpinolaRed]{PROWESS — Corpo}
  \azione{Finezza}{Finesse}{2}{Coltello, movimento silenzioso, velocità}
  \azione{Ombra}{Prowl}{3}{Pedinamento, appostamento, sparire in una folla}
  \azione{Combattimento}{Skirmish}{4}{Il lavoro che sa fare meglio. Decenni di pratica.}
  \azione{Forza bruta}{Wreck}{2}{Sfonda, intimidisce fisicamente}
\end{attributoblock}


\begin{attributoblock}[SpinolaRed]{RESOLVE — Spirito}
      \azione{Percezione}{Attune}{1}{Istinto per l'agguato e la trappola}
    \azione{Comando}{Command}{2}{Autorità fisica — la gente obbedisce quando lo guarda}
    \azione{Relazioni}{Consort}{1}{Bettole, mercenari, gente di strada}
    \azione{Persuasione}{Sway}{2}{Convince nel modo sbagliato ma efficace}
  \end{attributoblock}

\medskip
\begin{pgclockbox}{Stress \hfill \clock{9}{0}}
  \small\textit{Traumi: $\square$ Freddo \quad $\square$ Duro \quad $\square$ Ossessionato \quad $\square$ Paranoico \quad $\square$ Irrazionale \quad $\square$ Spezzato}
\end{pgclockbox}

\section*{Abilità Speciale}

% @rpg-schema: <.../pg-ranieri-costa>
%   fitd:specialAbility [ rdfs:label "Lavoro Pulito" ] .
\begin{spinolabox}{Lavoro Pulito}
  Quando Ranieri elimina, cattura o neutralizza un bersaglio, può scegliere
  di farlo senza lasciare tracce attribuibili agli Spinola. L'azione non
  genera \textbf{Calore} per la famiglia.

  \medskip\noindent\textit{Non richiede tiro aggiuntivo. Se il bersaglio è un
  PNG significativo (non una guardia generica), il GM può richiedere un tiro
  su \textsc{Ombra} per l'esecuzione silenziosa.}
\end{spinolabox}

\section*{Legami}

\begin{table}[h]
\centering\renewcommand{\arraystretch}{1.4}
\begin{tabularx}{\linewidth}{@{}L{2.6cm} X@{}}
  \toprule\rowcolor{SpinolaRed}
  \textcolor{white}{\textbf{PG}} & \textcolor{white}{\textbf{Legame}} \\
  \midrule
  Raffaele & Mi ha preso quando ero nessuno. Gli sono fedele. Ma fedeltà non significa cecità. \\
  \rowcolor{SpinolaRed!8}
  Margherita & L'agente di campo. Lavoriamo sullo stesso terreno con metodi diversi. Ci rispettiamo. \\
  Ippolito & Il ragazzo vuole giocare al soldato. Lo tengo d'occhio. Non voglio che si faccia del male. \\
  \rowcolor{SpinolaRed!8}
  Padre Anselmo & Il confessore. L'unico che sa quante cose ho fatto. Mi assolve ogni volta. Non capisco perché. \\
  \bottomrule
\end{tabularx}
\end{table}

\section*{Equipaggiamento}
\begin{goldlist}
  \item Spada corta e due coltelli — il terzo è nascosto nell'avambraccio sinistro
  \item Lista di cinque mercenari fidati che può radunare in un'ora nel porto
  \item Chiavi di quattro uscite secondarie di palazzi nel centro di Genova
  \item Borsellino separato con monete francesi — per le spese nei quartieri sotto controllo
\end{goldlist}


\begin{secretbox}
  \textbf{Segreto:} Tre mesi fa Ranieri ha ricevuto un incarico da Raffaele che
  ha eseguito senza fare domande. Ma dopo ha scoperto chi era la persona che aveva
  neutralizzato: un agente degli Spinola stessi, che stava per rivelare al
  governatore francese l'esistenza della memoria segreta preparata da Raffaele
  per Luigi XII. Ranieri ha quindi salvato Raffaele senza saperlo — e ora sa
  della memoria. Sa troppo.

  \medskip\noindent\textcolor{SecretRed}{\textbf{Agenda:}}
  \textit{Decidere cosa fare con quello che sa. Non è un ricattatore — ma non
  è nemmeno stupido. Se la rivolta va male e Raffaele consegna la famiglia
  ai francesi, Ranieri vuole avere qualcosa in mano per comprare la propria
  sopravvivenza.}
\end{secretbox}


% ════════════════════════════════════════════════════════════
% PG 7 — PADRE ANSELMO DA SAVONA
% @rpg-schema: <https://rpg-schema.org/instances/genova1507/pg-padre-anselmo>
%   a fitd:Character ;
%   rdfs:label "Padre Anselmo da Savona" ;
%   fitd:crew <.../crew-spinola> ;
%   fitd:role "Il Confessore / Il Ponte con la Chiesa" ;
%   fitd:action [ fitd:actionName "Study" ; fitd:dots 3 ] ;
%   fitd:action [ fitd:actionName "Attune" ; fitd:dots 3 ] ;
%   fitd:action [ fitd:actionName "Consort" ; fitd:dots 3 ] .
% ════════════════════════════════════════════════════════════

\clearpage

\pgheader[SpinolaRed]{Padre Anselmo da Savona}{Il Confessore / Il Ponte con la Chiesa}{Spinola}
\pgquote{La confessione è il sacramento della misericordia. È anche la fonte di informazioni più affidabile che esista.}

\section*{Background \& Motivazione}

\begin{rulebox}{}
  \textbf{Background:} Cinquantasei anni, agostiniano, confessore di famiglia degli
  Spinola da quindici anni. È diverso dai preti dei Fieschi — non predica, non
  mobilitizza folle, non ha ambizioni episcopali. Il suo potere è interamente
  privato e interstiziale: conosce i segreti di ogni membro della famiglia,
  ha accesso agli ambienti ecclesiastici che gli Spinola non possono frequentare
  direttamente, e ha sviluppato nel tempo una rete di confessori in tutta la
  Liguria che si scambiano informazioni con la discrezione professionale
  tipica della loro vocazione.

  \medskip\noindent\textcolor{SpinolaRed}{\textbf{Motivazione:}} La Chiesa non può
  servire bene i fedeli sotto un'occupazione straniera che non rispetta le
  immunità ecclesiastiche. La motivazione è sincera e strumentale allo stesso tempo.
  Padre Anselmo non vede contraddizione.
\end{rulebox}

\section*{Attributi \& Azioni}


\begin{attributoblock}[SpinolaRed]{INSIGHT — Mente}
  \azione{Caccia}{Hunt}{1}{Trova persone attraverso la rete parrocchiale}
  \azione{Studio}{Study}{3}{Teologia, diritto canonico, psicologia pratica}
  \azione{Osserva}{Survey}{2}{Legge le persone come testi — tra le righe}
  \azione{Ingegno}{Tinker}{1}{Messaggi cifrati, comunicazioni sicure via clero}
\end{attributoblock}

\columnbreak

\begin{attributoblock}[SpinolaRed]{PROWESS — Corpo}
  \azione{Finezza}{Finesse}{1}{Mani da scrivania, ma precise}
  \azione{Ombra}{Prowl}{1}{Si muove nelle sacrestie e nei corridoi ecclesiastici}
  \azione{Combattimento}{Skirmish}{0}{Mai}
  \azione{Forza bruta}{Wreck}{0}{No}
\end{attributoblock}


\begin{attributoblock}[SpinolaRed]{RESOLVE — Spirito}
      \azione{Percezione}{Attune}{3}{Sente il peso morale di una persona — e lo usa}
    \azione{Comando}{Command}{1}{Autorità spirituale, non fisica}
    \azione{Relazioni}{Consort}{3}{Clero basso e medio, confraternite, notai ecclesiastici}
    \azione{Persuasione}{Sway}{3}{Convince attraverso la colpa, la speranza e il debito morale}
  \end{attributoblock}

\medskip
\begin{pgclockbox}{Stress \hfill \clock{9}{0}}
  \small\textit{Traumi: $\square$ Freddo \quad $\square$ Duro \quad $\square$ Ossessionato \quad $\square$ Paranoico \quad $\square$ Irrazionale \quad $\square$ Spezzato}
\end{pgclockbox}

\section*{Abilità Speciale}

% @rpg-schema: <.../pg-padre-anselmo>
%   fitd:specialAbility [ rdfs:label "Rete Confessionale" ] .
\begin{spinolabox}{Rete Confessionale}
  Padre Anselmo può ottenere informazioni su qualsiasi persona che abbia
  un confessore cattolico — cioè praticamente chiunque a Genova nel 1507.
  Una volta per sessione, dichiara di conoscere un \textbf{segreto confessionale}
  di un PNG rilevante. L'informazione è accurata ma non può essere usata
  come prova pubblica.

  \medskip\noindent\textit{Non richiede tiro. L'informazione esiste.
  Usarla in modo pubblicamente attribuibile richiede un tiro su
  \textsc{Persuasione}.}
\end{spinolabox}

\section*{Legami}

\begin{table}[h]
\centering\renewcommand{\arraystretch}{1.4}
\begin{tabularx}{\linewidth}{@{}L{2.6cm} X@{}}
  \toprule\rowcolor{SpinolaRed}
  \textcolor{white}{\textbf{PG}} & \textcolor{white}{\textbf{Legame}} \\
  \midrule
  Raffaele & Lo conosco da quindici anni. Sa che so tutto. Io so che sa che so. Andiamo d'accordo. \\
  \rowcolor{SpinolaRed!8}
  Lucrezia & Si confessa raramente e sempre in modo evasivo. Questo mi dice più di quanto pensi. \\
  Ser Benedetto & Il notaio porta un peso che non dovrebbe portare da solo. Cerco di alleggerirglielo. Non ci riesco. \\
  \rowcolor{SpinolaRed!8}
  Ranieri & Il bravo. Lo assolvo perché me lo chiede e perché la misericordia non ha prerequisiti. \\
  \bottomrule
\end{tabularx}
\end{table}

\section*{Equipaggiamento}
\begin{goldlist}
  \item Breviario con un indice cifrato dei contatti nella rete confessionale ligure
  \item Lettera di credenziali del vescovo di Savona — apre porte in qualsiasi diocesi
  \item Registro personale in codice — trent'anni di osservazioni su famiglie e individui
  \item Piccola fiala di olio santo — e una di laudano, «per i casi di necessità pastorale»
\end{goldlist}


\begin{secretbox}
  \textbf{Segreto:} Padre Anselmo conosce sia il segreto di Raffaele (la memoria
  per i francesi) sia il segreto di Ranieri (che lo sa). È l'unico che ha il
  quadro completo. Non lo ha rivelato né all'uno né all'altro perché farlo
  distruggerebbe entrambe le relazioni su cui la famiglia Spinola poggia la
  propria capacità operativa. Porta questo sapere come un macigno.

  \medskip\noindent\textcolor{SecretRed}{\textbf{Agenda:}}
  \textit{Far sì che la rivolta abbia successo abbastanza in fretta da rendere
  irrilevante la memoria di Raffaele — e il segreto di Ranieri su di essa.
  Se il clock raggiunge 8, tutti si dimenticheranno di tutto. Questo è il
  suo piano. Non è un buon piano. Ma è l'unico che non richiede di tradire
  nessuno.}
\end{secretbox}

